\subsection{A theta-specific heat-flow truncation estimate}
\label{sec:heattail}
The coupled evolution equation alone does not determine the
sign of the Laguerre expressions. The square-lattice structure
nevertheless gives an error estimate uniform in the heat time
and in the real Fourier frequency. This separates the
approximation problem from the missing positivity argument.

Set $q_m=\pi m^2$ and
\[
H_{t,M}(z)=\int_0^\infty e^{tu^2}
\left(\sum_{m=1}^M
q_m(2q_me^{4u}-3)e^{5u-q_me^{4u}}\right)
\cos(zu)\,du .
\]
The complete function is $H_t$, while the
full-line transforms are $F_t=2H_t$ and $F_{t,M}=2H_{t,M}$.

\begin{theorem}[Uniform theta-tail bounds along the heat flow]
\label{thm:uniformheattail}
For all $M\ge1$, $k\ge0$,
$t\in[0,\tfrac12]$ and
$z\in\mathbb C$ with $|\Im z|\le3$,
put $q=\pi(M+1)^2$. Then
\begin{equation}
\left|F_t^{(k)}(z)-F_{t,M}^{(k)}(z)\right|
\le\frac{k!}{4^k}\,3q\,e^{-q+1/16}.
\label{eq:uniformheattail}
\end{equation}
For $k=0$ the same estimate is valid
on the larger strip $|\Im z|\le7$.
These bounds are simultaneous in $t$ and
in $\Re z$, and require no assumption
about the zeros of $H_t$.
\end{theorem}
\begin{proof}
For $u\ge0$ the elementary exponential
Taylor inequality $e^{4u}\ge8u^2$
implies $tu^2\le e^{4u}/16$.
Also $u^k\le(k!/4^k)e^{4u}$ and
$|\sin(zu)|,|\cos(zu)|\le e^{|\Im z|u}$.
Since $q_m>3$ and each theta
summand is positive for $u\ge0$,
its absolute value is bounded by
$2q_m^2e^{9u-q_me^{4u}}$.
For $a_m=q_m-1/16$ and $|\Im z|\le3$
we obtain
\begin{align*}
&\left|F_t^{(k)}(z)-F_{t,M}^{(k)}(z)\right|\\
&\le\frac{k!}{4^k}
\sum_{m>M}q_m^2\int_1^\infty
v^3e^{-a_mv}\,dv\\
&=\frac{k!}{4^k}\sum_{m>M}
q_m^2e^{-a_m}
\left(\frac1{a_m}+\frac3{a_m^2}
+\frac6{a_m^3}+\frac6{a_m^4}\right).
\end{align*}
For $q_m\ge12$, rational algebra gives
\begin{equation}
2q_m(a_m^3+3a_m^2+6a_m+6)
<3a_m^4 .
\label{eq:heatrationalgate}
\end{equation}
Indeed the difference between the right
and left sides, with $q_m=12+v$, is
\[
v^4+\frac{333}8v^3+
\frac{39891}{64}v^2+
\frac{7941835}{2048}v+
\frac{522416259}{65536}>0 .
\]
Thus every term is bounded by
$\tfrac32q_me^{-q_m+1/16}$.
For $m\ge2$ the ratios of successive
$q_me^{-q_m}$ are bounded by
$4e^{-9}<1/2$. Summing the geometric
majorant gives \eqref{eq:uniformheattail}.
When $k=0$ the auxiliary $e^{4u}$
is unnecessary; the exponent of $v$
is $(5+|\Im z|)/4\le3$ throughout
$|\Im z|\le7$, proving the second claim.
\end{proof}

\begin{corollary}[Effective approximation of the coupled sources]
\label{cor:heatsources}
Let $L_{n,t}$ and $L_{n,t,M}$ denote
the generalized Laguerre expressions of
$F_t$ and $F_{t,M}$.
Define the finite positive majorant
\[
C_*=\sum_{m\ge1}q_m^2e^{-(q_m-1/16)}
\left(\frac1{a_m}+\frac3{a_m^2}
+\frac6{a_m^3}+\frac6{a_m^4}\right).
\]
For every $n\ge0$,
$t\in[0,\tfrac12]$ and $|\Im z|\le3$,
\begin{equation}
|L_{n,t}(z)-L_{n,t,M}(z)|
\le \frac{6(2n+1)C_*}{4^{2n}}\,
q e^{-q+1/16}.
\label{eq:heatlagerror}
\end{equation}
Consequently the coupling source
$S_{n,t}=(n+1)(2n+1)L_{n+1,t}$
has the same explicit bound,
with $n$ replaced by $n+1$ and
multiplied by $(n+1)(2n+1)$.
\end{corollary}
\begin{proof}
The estimate in the preceding theorem,
applied to the complete sum, bounds
$|F_t^{(j)}(z)|$ and
$|F_{t,M}^{(j)}(z)|$ by
$C_*j!/4^j$.
Expand the difference of each product
in the finite Laguerre derivative sum.
The identity
$\binom{2n}{j}j!(2n-j)!=(2n)!$
cancels the factorials, leaving
$2n+1$ identical upper bounds and
the constant in \eqref{eq:heatlagerror}.
The source estimate follows
immediately from its definition.
\end{proof}

\paragraph{What the estimate does not give.}
Theorem~\ref{thm:uniformheattail} controls
the actual theta heat family, simultaneously
for all real frequencies and heat times up
to $1/2$. It is therefore a usable input
to a certified finite-height contour
analysis and to any proposed estimate of
the adjacent-order heat coupling.
It is an \emph{absolute} error estimate.
At zeros or near zeros of $L_{n,t}$
it gives no relative control over
$L_{n+1,t}/L_{n,t}$ and cannot imply
the comparison hypothesis
\eqref{eq:comparisonmissing}.
The precise missing assertion is a
global-in-frequency, all-order source
control or a different theta-specific
complex zero-exclusion mechanism.
No such assertion is proved here.

The polynomial comparison
\eqref{eq:heatrationalgate} and the
factorial cancellation are checked in
\path{XiKernel/ThetaHeatTail.lean}
and independently by
\path{certificates/xi_theta_heat_tail_exact.py}.
The infinite theta-series, exponential
Taylor domination and Fourier-integral
transfer remain analytic arguments,
not end-to-end Lean theorems.
